\chapter{Setting: betting wealth processes}
\label{chap:setting}

We test the null hypothesis $H_0 : \E_P[X] = m$ on i.i.d.\ observations $X_0, X_1, \dots$ with
values in $[0, 1]$ and law $P$. A distribution $P$ on $[0,1]$ with mean $m$ is a \emph{null
distribution}; it is \emph{non-degenerate} if $P \ne \delta_m$.

\section{Wealth processes}

\begin{definition}[Range of the bet fractions]
  \label{def:fraction_range}
  \lean{Learning.Betting.fractionRange}
  \leanok
  $\Fr := [-\frac{1}{1-m}, \frac{1}{m}]$. For $m \in (0,1)$, $x \in [0,1]$ and $\lambda \in \Fr$,
  $1 + \lambda(x - m) \ge 0$.
\end{definition}

\begin{lemma}[Basic properties of the range]
  \label{lem:fraction_range_basic}
  \lean{Learning.Betting.isCompact_fractionRange, Learning.Betting.interior_fractionRange,
    Learning.Betting.zero_mem_fractionRange, Learning.Betting.mem_fractionRange_of_mem_uIcc,
    Learning.Betting.Icc_neg_one_one_subset_interior_fractionRange,
    Learning.Betting.Icc_neg_one_one_subset_fractionRange,
    Learning.Betting.abs_le_of_mem_fractionRange}
  \leanok
  \uses{def:fraction_range}
  $\Fr$ is a compact interval with interior $(-\frac{1}{1-m}, \frac{1}{m})$, containing $0$ and
  hence every $\lambda'$ between $0$ and any of its points; every $\lambda \in \Fr$ satisfies
  $\abs\lambda \le \max(\frac1m, \frac{1}{1-m})$. For $m \in (0,1)$, $[-1, 1]$ is contained in the
  interior of $\Fr$.
\end{lemma}

\begin{proof}
  \leanok
  Immediate; $[-1,1] \subseteq (-\frac{1}{1-m}, \frac1m)$ since $\frac1m, \frac{1}{1-m} > 1$.
\end{proof}

With the convention $1/0 = 0$ of Lean, these statements, and the nonnegativity of the wealth below,
hold for $m \in [0, 1]$, which is the range given by the null hypothesis.


\begin{definition}[Wealth]
  \label{def:wealth}
  \lean{Learning.Betting.wealth}
  \leanok
  For real processes $\lambda = (\lambda_n)$ (the bet fractions) and $X = (X_n)$ (the observations),
  the wealth after $n$ rounds is $W^\lambda_n := \prod_{k < n} (1 + \lambda_k (X_k - m))$.
\end{definition}

\begin{lemma}[Basic properties of the wealth]
  \label{lem:wealth_basic}
  \lean{Learning.Betting.wealth_zero, Learning.Betting.wealth_succ, Learning.Betting.wealth_nonneg,
    Learning.Betting.measurable_wealth, Learning.Betting.one_add_mul_sub_nonneg}
  \leanok
  \uses{def:wealth, def:fraction_range}
  $W_0 = 1$, $W_{n+1} = W_n (1 + \lambda_n (X_n - m))$; if $m \in [0,1]$, $\lambda_k \in \Fr$ and
  $X_k \in [0,1]$ for $k < n$ then $W_n \ge 0$; $W_n$ is measurable when the $\lambda_k$ and $X_k$
  are.
\end{lemma}

\begin{proof}
  \leanok
  Immediate from the definition; $1 + \lambda(x - m) = (1-x)(1 - \lambda m) + x(1 + \lambda(1-m))$ is
  a convex combination of nonnegative numbers.
\end{proof}

\begin{definition}[Fixed-fraction, mixture and hedged wealths]
  \label{def:fixed_wealth}
  \lean{Learning.Betting.fixedWealth, Learning.Betting.mixtureWealth, Learning.Betting.hedgedWealth}
  \leanok
  \uses{def:wealth}
  The fixed-fraction wealth $W^\lambda_n$ for a constant $\lambda \in \R$; the mixture wealth
  $W^\pi_n := \int W^\lambda_n \, d\pi(\lambda)$ for a probability measure $\pi$ on $\Fr$; the hedged
  wealth $W^{\pm\lambda}_n := \frac12 (W^\lambda_n + W^{-\lambda}_n)$ for a process $\lambda$.
\end{definition}

\begin{definition}[Best-in-hindsight log-wealth and regret]
  \label{def:hindsight}
  \lean{Learning.Betting.hindsightLogWealth, Learning.Betting.logRegret}
  \leanok
  \uses{def:fixed_wealth}
  $L^*_n := \log \sup_{\lambda \in \Fr} W^\lambda_n$ and, for a strategy $\lambda$,
  $R_n := L^*_n - \log W^\lambda_n$.
\end{definition}

\begin{remark}
  The paper defines $L^*_n = \max_{\lambda \in \Fr} \log W^\lambda_n$, which is the same number
  when $W^\lambda_n > 0$ for some $\lambda$ (e.g.\ $\lambda = 0$). We take the logarithm of the
  supremum to avoid Lean's convention $\log 0 = 0$ at the bet fractions with zero wealth. By the
  dual representation of \cite{honda2010asymptotically}, $L^*_n = n \klinf(P_n, m)$ for the
  empirical measure $P_n$; this identity is not used.
\end{remark}

\begin{definition}[Sub-Gaussian test processes]
  \label{def:subg_test}
  \lean{Learning.Betting.subgaussianTest, Learning.Betting.subgaussianMixtureTest}
  \leanok
  For a real process $\lambda$, $M^\lambda_n := \exp\big(\sum_{k<n} \frac{(X_k - m)^2 - (X_k -
  \lambda_k)^2}{2}\big)$, and for a probability measure $\pi$ on $\R$,
  $M^\pi_n := \int M^\lambda_n \, d\pi(\lambda)$ (constant $\lambda$ in the integrand).
\end{definition}

\section{Betting strategies}

\begin{definition}[Predictable strategies]
  \label{def:predictable}
  \leanok
  Let $\F$ be a filtration of $(\Omega, \mathcal A)$. The observations are \emph{i.i.d.\ along
  $\F$} with law $P$ if each $X_n$ is $\F_{n+1}$-measurable, independent of $\F_n$, and has law
  $P$. A \emph{predictable strategy} is a process $\lambda$ adapted to $\F$ ($\lambda_n$ is
  $\F_n$-measurable). In Lean these are the hypotheses
  \texttt{hX}, \texttt{hindep}, \texttt{hlaw}, \texttt{hlam} of the filtration forms of the results
  (e.g.\ \texttt{Learning.Betting.sum\_sq\_criterion\_of\_indep}).
\end{definition}

\begin{definition}[The i.i.d.\ environment]
  \label{def:iid_env}
  \lean{Learning.Environment.const}
  \leanok
  For a probability measure $P$ on $\R$, $\mathrm{const}(P)$ is the LML environment (with trivial
  observations) whose feedback is drawn from $P$ at each round, whatever the action. A run of an
  algorithm $\mathrm{alg}$ with actions in $\R$ in this environment is a pair of processes
  $(\lambda, X)$ such that $\lambda_n$ is drawn by $\mathrm{alg}$ (possibly at random) from the
  history of the first $n$ rounds and $X_n \sim P$ is independent of that history and of
  $\lambda_n$.
\end{definition}

\begin{lemma}
  \label{lem:iid_env_feedback}
  \lean{Learning.IsAlgEnvSeq.hasLaw_feedback_const, Learning.IsAlgEnvSeq.indepFun_feedback_const,
    Learning.IsAlgEnvSeq.iIndepFun_feedback_const}
  \leanok
  \uses{def:iid_env}
  In a run of an algorithm in $\mathrm{const}(P)$, each $X_n$ has law $P$, is independent of the
  history of the first $n$ rounds and of $\lambda_n$, and the $X_n$ are mutually independent.
\end{lemma}

\begin{proof}
  \leanok
  The conditional distribution of $X_n$ given the history and $\lambda_n$ is the constant kernel
  $P$.
\end{proof}

\begin{lemma}
  \label{lem:iid_of_run}
  \lean{Learning.IsAlgEnvSeq.adapted_action_filtrationAction,
    Learning.IsAlgEnvSeq.measurable_feedback_filtrationAction_succ,
    Learning.IsAlgEnvSeq.indep_feedback_filtrationAction_const}
  \leanok
  \uses{def:iid_env, def:predictable, lem:iid_env_feedback}
  In a run $(\lambda, X)$ of an algorithm in $\mathrm{const}(P)$, the observations are i.i.d.\ along
  the filtration generated by the history of the first $n$ rounds and the action $\lambda_n$
  (LML's \texttt{IsAlgEnvSeq.filtrationAction}), and $\lambda$ is adapted to it.
\end{lemma}

\begin{proof}
  \leanok
  $X_n$ is a function of the history of the first $n + 1$ rounds, $\lambda_n$ of the generator of
  the $n$-th $\sigma$-algebra, and Lemma~\ref{lem:iid_env_feedback} gives the independence and the
  law.
\end{proof}

\begin{definition}[Plug-in strategies]
  \label{def:plug_in}
  \lean{Learning.Betting.IsPlugIn}
  \leanok
  $\lambda$ is a \emph{plug-in} strategy for $X$ if each $\lambda_n$ is a measurable function of
  $(X_0, \dots, X_{n-1})$.
\end{definition}

\begin{definition}[Past filtration]
  \label{def:past_filtration}
  \lean{Learning.Betting.pastFiltration}
  \leanok
  The \emph{past filtration} of a process $X$ is $\F^X_n := \sigma(X_0, \dots, X_{n-1})$.
\end{definition}

\begin{lemma}
  \label{lem:iid_of_plug_in}
  \lean{Learning.Betting.measurable_pastFiltration_succ, Learning.Betting.indep_pastFiltration,
    Learning.Betting.IsPlugIn.adapted_pastFiltration}
  \leanok
  \uses{def:past_filtration, def:plug_in, def:predictable}
  If the $X_n$ are independent with law $P$ (and measurable), they are i.i.d.\ along $\F^X$, and
  every plug-in strategy for $X$ is adapted to $\F^X$.
\end{lemma}

\begin{proof}
  \leanok
  $X_n$ is independent of $(X_0, \dots, X_{n-1})$ by mutual independence (disjoint index sets).
\end{proof}

\begin{definition}[Classical strategies]
  \label{def:strategies}
  \lean{Learning.Betting.ktFraction, Learning.Betting.IsGrapa, Learning.Betting.empMean,
    Learning.Betting.empVar, Learning.Betting.agrapaFraction, Learning.Betting.hedgeFraction}
  \leanok
  \uses{def:fixed_wealth, def:plug_in}
  With $S_n := \sum_{k<n} (X_k - m)$, $\hat\mu_n := \frac1n \sum_{k<n} X_k$ and
  $\hat\sigma^2_n := \frac1n \sum_{k<n} (X_k - \hat\mu_n)^2$:
  \begin{itemize}
    \item Krichevsky--Trofimov: $\lKT_n := \frac{1/2 + S_n}{C (n+1)}$;
    \item GRAPA: $\lgr$ is a plug-in strategy such that $\lgr_n \in \Fr$ maximizes
          $\lambda \mapsto W^\lambda_n$ over $\Fr$ (the paper's $\lambda^{\mathsf{GRAPA}}_{n+1} =
          \lambda^{\mathsf{KL}}_n$);
    \item approximate GRAPA: $\lagr_n := \big(-\frac{C}{1-m}\big) \vee \frac{\hat\mu_n -
          m}{\hat\sigma^2_n + (\hat\mu_n - m)^2} \wedge \frac{C}{m}$;
    \item predictable hedging with variance estimate $v$: $\lPrH_n := \big(-\frac{C}{1-m}\big)
          \vee \sqrt{\frac{2 \log(2/\alpha)}{v_n (n+1) \log(n+2)}} \wedge \frac{C}{m}$.
  \end{itemize}
\end{definition}

\begin{lemma}
  \label{lem:isPlugIn_kt}
  \lean{Learning.Betting.isPlugIn_ktFraction}
  \leanok
  \uses{def:strategies}
  The KT strategy is a plug-in strategy.
\end{lemma}

\begin{proof}
  \leanok
  $\lKT_n$ is a continuous function of $X_0, \dots, X_{n-1}$.
\end{proof}

\section{Asymptotic notation in probability}

\begin{definition}[Landau notation in probability and almost surely]
  \label{def:asymp}
  \lean{ProbabilityTheory.IsBigOInProb, ProbabilityTheory.IsBigOmegaInProb,
    ProbabilityTheory.IsBigOAE, ProbabilityTheory.IsBigOmegaAE}
  \leanok
  For random variables $Z_n$ and numbers $a_n$: $Z_n = \Op(a_n)$ if for every $\eps > 0$ there is
  $M > 0$ with $\Prob(\abs{Z_n} \le M a_n) \ge 1 - \eps$ for all but finitely many $n$;
  $Z_n = \Omp(a_n)$ if for every $\eps > 0$ there is $\delta > 0$ with
  $\Prob(\abs{Z_n} \ge \delta a_n) \ge 1 - \eps$ for all but finitely many $n$; $Z_n = \Oas(a_n)$
  (resp.\ $\Omas(a_n)$) if almost surely $Z_n = O(a_n)$ (resp.\ $a_n = O(Z_n)$).
\end{definition}

\section{Leverage and portfolios}

\begin{definition}[Minimal next wealth, predictably non-bankrupt event]
  \label{def:pnb}
  \lean{Wang2026Almost.nextWealthMin, Wang2026Almost.pnbEvent}
  \leanok
  \uses{def:wealth}
  $\Wh_n := \min_{x \in \{0, 1\}} W_n (1 + \lambda_n (x - m))$, the minimal wealth after round $n$
  given the information available before $X_n$; for $\rho \in \R$, the \emph{$\rho$-predictably
  non-bankrupt event} is $N^\rho := \bigcap_n \{\Wh_n > \rho\}$.
\end{definition}

\begin{definition}[Leverage]
  \label{def:leverage}
  \lean{Wang2026Almost.leverageFraction, Wang2026Almost.oppLeverage}
  \leanok
  \uses{def:pnb}
  The leveraged fraction $\beta^\rho_n := \frac{W_n \lambda_n}{W_n - \rho}$ (borrowing
  $\frac{\rho}{1-\rho}$ units of cash and betting $\lambda$ with the leveraged wealth), and the
  \emph{$\rho$-opportunistic leverage} $\gamma_n := \beta^\rho_n$ if $\Wh_n > \rho$, $\gamma_n :=
  \lambda_n$ otherwise.
\end{definition}

\begin{definition}[Portfolio of two strategies]
  \label{def:portfolio}
  \lean{Wang2026Almost.portfolioFraction}
  \leanok
  \uses{def:wealth}
  For strategies $\lambda, \nu$ and $\kappa \in \R$,
  $\beta_n := \frac{(1-\kappa) W^\lambda_n \lambda_n + \kappa W^\nu_n \nu_n}{(1-\kappa)W^\lambda_n +
  \kappa W^\nu_n}$.
\end{definition}

\begin{definition}[Paths]
  \label{def:paths}
  \lean{Wang2026Almost.coord, Wang2026Almost.IsBinaryPath, Wang2026Almost.IsPathPredictable}
  \leanok
  On the path space $\R^{\N}$: the coordinate process, the binary paths (values in $\{0,1\}$),
  and the strategies $\lambda_n(x)$ depending only on $x_0, \dots, x_{n-1}$.
\end{definition}

\section{Two mixing distributions}

\begin{definition}[Universal portfolio and Robbins' mixture]
  \label{def:mixtures}
  \lean{Wang2026Almost.betaMixture, Wang2026Almost.robbinsConst, Wang2026Almost.robbinsDensity,
    Wang2026Almost.robbinsMixture}
  \leanok
  $\pi_{a,b}$ is the image of the $\mathrm{Beta}(a, b)$ distribution by $x \mapsto 2x - 1$;
  Robbins' mixing distribution has density
  $f(\lambda) = \frac{\indic\{\abs{\lambda} \le 1\} \log\log C}{2 \abs\lambda \log(C/\abs\lambda)
  (\log\log(C/\abs\lambda))^2}$ with respect to the Lebesgue measure, $C = 6.6 e$.
\end{definition}

\section{Basic consequences of the null hypothesis}

\begin{lemma}
  \label{lem:null_basic}
  \lean{Learning.Betting.mem_Icc_of_integral_eq, Learning.Betting.memLp_id_of_mem_Icc,
    Learning.Betting.mem_Ioo_of_ne_dirac, Learning.Betting.variance_pos_of_ne_dirac,
    Learning.Betting.mem_Ioo_of_variance_pos, Learning.Betting.ne_dirac_of_variance_pos,
    MeasureTheory.Measure.eq_dirac_of_ae_eq, ProbabilityTheory.HasLaw.ae_comp}
  \leanok
  Let $P$ be a distribution on $[0,1]$ with mean $m$. Then $m \in [0,1]$, $P$ has moments of all
  orders and, if $X \sim P$, $X \in [0,1]$ almost surely. If $P$ is non-degenerate, then
  $m \in (0, 1)$ and $\sigma^2 := \Var_P > 0$; conversely $\sigma^2 > 0$ implies that $P$ is
  non-degenerate.
\end{lemma}

\begin{proof}
  \leanok
  If $\Var_P = 0$ then $P = \delta_m$. If $m = 0$ then $X = 0$ a.s.\ ($X \ge 0$ with mean $0$), so
  $P = \delta_0$; similarly for $m = 1$.
\end{proof}

\begin{lemma}
  \label{lem:null_subgaussian}
  \lean{Learning.Betting.hasSubgaussianMGF_sub_of_mem_Icc,
    Learning.Betting.hasSubgaussianMGF_sq_sub_variance, ProbabilityTheory.HasLaw.hasSubgaussianMGF_comp}
  \leanok
  \uses{lem:hoeffding_lemma_null}
  Let $P$ be a distribution on $[0,1]$ with mean $m$ and variance $\sigma^2$. Under $P$, $x - m$ and
  $(x - m)^2 - \sigma^2$ are $\frac14$-sub-Gaussian; so are $X - m$ and $(X - m)^2 - \sigma^2$ for
  $X \sim P$.
\end{lemma}

\begin{proof}
  \leanok
  Lemma~\ref{lem:hoeffding_lemma_null}: both are centered with values in an interval of length $1$
  ($[-m, 1-m]$ and $[-\sigma^2, 1 - \sigma^2]$, as $(x - m)^2 \in [0, 1]$).
\end{proof}
